La distribución de los rendimientos diarios confirma diferencias importantes en la volatilidad entre las empresas analizadas.
TSLA
,
PDD
y
BNTX
presentan las distribuciones más amplias, lo que indica una mayor dispersión de los rendimientos diarios y, por tanto, un comportamiento más volátil. Este resultado es consistente con los análisis previos, donde estas compañías también mostraron mayores fluctuaciones en los rendimientos acumulados y mensuales.
Por el contrario,
SNPS
,
KLAC
,
FICO
,
CDNS
y
ARES
exhiben distribuciones considerablemente más concentradas alrededor del rendimiento promedio, reflejando una menor dispersión y una mayor estabilidad en las variaciones diarias de sus precios. Por su parte,
AMD
y
NVDA
ocupan una posición intermedia, mostrando una dispersión superior a la de las empresas más estables, pero inferior a la observada en
TSLA
,
PDD
y
BNTX
.
Finalmente, se observa que todas las distribuciones se encuentran centradas alrededor del
0 %
, lo que indica que las diferencias entre empresas no radican en el rendimiento diario promedio, sino en la magnitud de las variaciones que experimentan día a día. En consecuencia, las empresas con distribuciones más amplias implican un mayor nivel de riesgo, ya que presentan cambios más pronunciados en el precio de sus acciones.